\subsection{Hierarchical Nesting and PWC}

The prime-factor structure algorithm (ApportionRegions) builds districts
through a hierarchical nesting that follows the prime factorisation of $k$.
For NC ($k = 14 = 2 \times 7$), the algorithm first splits the state into
2 super-regions, then each super-region into 7 districts.

\paragraph{Why nesting reduces PWC.}
At each level of the hierarchy, the bisection split creates sub-regions that
are geographically compact \emph{within the parent region}.
Population concentrations (urban centres) that are geographically close
tend to be assigned to the same sub-region at each level, because METIS
minimises cut weight (edge crossings), which is higher for cuts that separate
spatially adjacent high-population nodes.

The result is that prime-factor districts tend to have their population
concentrated near the district centroid rather than dispersed to the
district periphery---a property that directly minimises PWC.

\paragraph{Standard-bisect comparison.}
Standard-bisect applies METIS to the full state graph and splits into $k$ parts
simultaneously.
Without hierarchical nesting, the algorithm may assign geographically distant
tracts to the same district (across multiple METIS cut levels), producing districts
with higher PWC.

\subsection{Quantitative Evidence}

The 14\% PWC reduction from prime-factor vs.\ standard-bisect on NC 2020
is computed as:
\[
  \text{reduction} = \frac{\text{PWC}_{\text{std}} - \text{PWC}_{\text{pf}}}{\text{PWC}_{\text{std}}}
  = \frac{72.1 - 62.1}{72.1} \approx 0.139 \approx 14\%.
\]

\subsection{Trade-off with Geometric Compactness}

Prime-factor does not lead on geometric compactness metrics.
Ratio-optimal achieves higher PP and Reock than prime-factor (K.1, K.2).
This trade-off---prime-factor minimises PWC at the cost of slightly lower geometric
compactness---reflects the orthogonality of representational and geometric objectives.
Practitioners who prioritise PWC for VRA analysis or community-of-interest
arguments should use prime-factor; those who prioritise PP for litigation
should use ratio-optimal.
